% !TEX program = xelatex
%
% CV of Sebastian Koch — LaTeX port of the MS Word template cv_2026_3page.docx.
% Compile with XeLaTeX (fonts: Roboto, Arial, Times New Roman via fontspec;
% Carlito from TeX Live stands in for the URL font in the header).
\documentclass{article}

\usepackage[letterpaper,
            left=0.5in, right=0.5in,
            top=0.591in, bottom=0.7in,
            footskip=12.5pt]{geometry}
\usepackage{fontspec}
\usepackage{xcolor}
\usepackage{longtable}
\usepackage{array}
\usepackage{graphicx}
\usepackage{fontawesome5}
\usepackage{fancyhdr}
\usepackage[hidelinks]{hyperref}

% ---------------------------------------------------------------- fonts ----
\setmainfont{Roboto}
\newfontfamily\headerserif{Times New Roman}
\newfontfamily\headersans{Arial}
\newfontfamily\urlfont{Carlito-Regular.ttf}[
  BoldFont=Carlito-Bold.ttf, ItalicFont=Carlito-Italic.ttf]

% Word line grid (measured from the reference PDF at 150 dpi):
%   10.5pt body on a 13.49pt grid, bullets add 0.96pt paragraph spacing,
%   compact sections (honors/publications/talks) use a 12.53pt grid.
\newcommand{\bodysize}{\fontsize{10.5}{13.49}\selectfont}
\newcommand{\compactsize}{\fontsize{10.5}{12.53}\selectfont}
\newcommand{\datesize}{\fontsize{9}{11.56}\selectfont}
\newcommand{\bulletgap}{\vskip 0.96pt}

% -------------------------------------------------------------- colours ----
\definecolor{dategray}{HTML}{767171}
\definecolor{footgray}{HTML}{7F7F7F}

% -------------------------------------------------------------- layout -----
% Word table geometry: date column, 0.075in padding on each side of the
% 1pt gray divider, content column filling the rest of the line.
\newlength{\dcw}   % date column width
\newlength{\ccw}   % content column width
\newcommand{\datebar}{\kern0.075in{\color{dategray}\vrule width 1pt}\kern0.075in}
\newcommand{\datestyle}{\raggedleft\arraybackslash\datesize\leavevmode\color{dategray}}

% \begin{cvtable}[<content font>]{<date column width>}
\newenvironment{cvtable}[2][\bodysize]{%
  \setlength{\dcw}{#2}%
  \setlength{\ccw}{\dimexpr\textwidth-\dcw-0.314in\relax}%
  \setlength{\LTpre}{3.2pt}\setlength{\LTpost}{0pt}%
  \setlength{\LTleft}{0pt}\setlength{\LTright}{\fill}%
  \setlength{\tabcolsep}{0pt}%
  \begin{longtable}{@{\kern0.075in}>{\datestyle}p{\dcw}!{\datebar}>{\raggedright\arraybackslash#1}p{\ccw}@{}}%
}{\end{longtable}\vspace{\entrygap}}

% Company line in the Experience section: full-width row without the divider.
\newcommand{\company}[1]{\multicolumn{2}{@{}l@{}}{\kern0.194in\textbf{#1}}\\}

% Word's Symbol-font bullet is much larger than Roboto's \textbullet.
\newcommand{\cvbullet}{\raisebox{-0.185em}{\scalebox{1.6}{\textbullet}}}

% Bullet items: bullet hangs left of the text, text indented 0.42in.
% \bitem is left-aligned (Word default), \bitemj justified (Word jc=both),
% \bitems is the small 9.5pt variant.
\newcommand{\bitem}[1]{\bulletgap{\setlength{\leftskip}{0.4215in}\noindent
  \llap{\cvbullet\kern0.12in}#1\par}}
\newcommand{\bitemj}[1]{\bulletgap{\setlength{\leftskip}{0.4215in}%
  \rightskip=0pt \parfillskip=0pt plus 1fil \noindent
  \llap{\cvbullet\kern0.12in}#1\par}}
\newcommand{\bitems}[1]{\bulletgap{\fontsize{9.5}{12.2}\selectfont
  \setlength{\leftskip}{0.4215in}\noindent\llap{\cvbullet\kern0.12in}#1\par}}

% Continuation bullet: no leading glue (for cells that start a table row).
\newcommand{\bitemcont}[1]{{\setlength{\leftskip}{0.4215in}%
  \rightskip=0pt \parfillskip=0pt plus 1fil \noindent
  \llap{\cvbullet\kern0.12in}#1\par}}

% Publication entry: hanging indent of 0.221in (Word ind l=318, hanging=318).
\newcommand{\pubitem}[1]{{\hangindent=0.221in \hangafter=1 \noindent#1\par}}

% ------------------------------------------------------------ sections -----
\newlength{\entrygap}\setlength{\entrygap}{4.6pt}
% \cvsection[<space above>]{<title>}
\newcommand{\cvsection}[2][14pt]{%
  \par\vspace{#1}%
  {\noindent\fontsize{11}{13.2}\selectfont\bfseries #2\par}%
  \vspace{2pt}\hrule height 1pt
  \vspace{5pt}}

% -------------------------------------------------------------- footer -----
\pagestyle{fancy}
\fancyhf{}
\renewcommand{\headrulewidth}{0pt}
\fancyfoot[R]{\color{footgray}\fontsize{9}{10.8}\selectfont Updated July 2026}

\hyphenpenalty=10000 % no auto-hyphenation (Word); breaks at explicit hyphens stay allowed
\setlength{\parindent}{0pt}
\setlength{\parskip}{0pt}

\begin{document}\bodysize

% ================================================================ header ====
% The Word template pulls the header table above the top margin.
\vspace*{-40pt}
\begin{center}
  {\headersans\fontsize{24}{26}\selectfont S{\fontsize{18.7}{26}\selectfont EBASTIAN}
   K{\fontsize{18.7}{26}\selectfont OCH}}\\[2pt]
  {\headerserif\itshape\fontsize{11}{13.2}\selectfont Research Scientist at Google}
\end{center}
\vspace{-7.7pt}
\noindent
\makebox[0.3333\textwidth][c]{\headerserif\fontsize{11}{13.2}\selectfont
  \raisebox{-1pt}{\faEnvelope}\kern0.55em
  \href{mailto:koch.sebastian@mail.com}{koch.sebastian@mail.com}}%
\makebox[0.3333\textwidth][c]{%
  \raisebox{-1pt}{\faGlobe}\kern0.4em
  {\urlfont\fontsize{10.5}{12.6}\selectfont
   \href{https://kochsebastian.com}{https://kochsebastian.com}}}%
\makebox[0.3333\textwidth][c]{\headerserif\fontsize{11}{13.2}\selectfont
  \raisebox{-1.5pt}{\faAddressCard}\kern0.55em Munich, Germany}
\par\vspace{2.9pt}
\hrule height 0.75pt

% ============================================================= education ====
\cvsection[15pt]{Education}

\begin{cvtable}{0.831in}
Apr 2022 --\par Dec 2026* &
  \textbf{Ulm University, Germany}\par
  PhD in Computer Science\par
  \bitem{\textit{Advisor}: \href{https://viscom.uni-ulm.de/members/timo-ropinski/}{Prof.\ Timo Ropinski}}
  \bitem{\textit{Thesis}: Language-Grounded Representations for Compositional 3D Scene Understanding}
  \bitem{\textit{Apr 2022 -- Apr 2025}: Bosch Research Doctoral Program}
  \bitem{\textit{May 2025 -- Nov 2025}: Google Student Researcher}
  \bitem{\textit{Mar 2026 -- Jul 2026}: Google Research Intern}\\
\end{cvtable}

\begin{cvtable}{0.831in}
2020 -- 2022 &
  \textbf{Tübingen University, Germany}\par
  Master of Science in Computer Science\par
  \bitem{\textit{Advisors}: Prof.\ \href{https://alr.anthropomatik.kit.edu/21_65.php}{Gerhard Neumann} \& \href{https://www.cvlibs.net/}{Prof.\ Andreas Geiger}}
  \bitem{\textit{Thesis project}: Multi-View RGB-D Fusion for 6D Pose Estimation}
  \bitem{\textit{GPA}: 1.4 (Scale: 1.0 [best] to 5.0 [fail])}
  \bitem{\textit{Selected courses}: Computer Vision, Deep Learning, 3D Computer Vision, Mobile Robotics, Reinforcement Learning, Mathematics for Machine Learning}\\
\end{cvtable}

\begin{cvtable}{0.831in}
2016 -- 2020 &
  \textbf{Baden-Württemberg Cooperative State University (DHBW), Germany}\par
  Bachelor of Engineering in Computer Science\par
  \bitem{\textit{Thesis project}: Improvement of the robustness of a SLAM system using Object Semantics}
  \bitem{\textit{GPA}: 1.8 (Scale: 1.0 [best] to 5.0 [fail])}\\
\end{cvtable}

% ============================================================ experience ====
\cvsection[10.2pt]{Experience}

\begin{cvtable}{0.831in}
\company{Google, Munich Germany}
Jul 2026 --\par Present &
  Research Scientist -- Immersive \& Spatial AI\par
  \bitemj{\textit{Role}: Research on approaches for Spatial and Immersive AI for Android XR, including Gemini, 3D Scene Graphs, Vision Foundation Models and more.}\\
\end{cvtable}

\begin{cvtable}{0.831in}
\company{Google Research, Zurich Switzerland}
Mar 2026 --\par Jul 2026 &
  Research Intern -- 3D Scene Reconstruction and Generation for Android XR\par
  \bitem{\textit{Advisors}: \href{https://scholar.google.com/citations?user=eQ0om98AAAAJ&hl=en}{Marie-Julie Rakotosaona} \& \href{https://federicotombari.github.io/}{Federico Tombari}}
  \bitemj{\textit{Role}: Research on the joint benefits of 3D reconstruction and world models with 3D Gaussian Splatting and video generation models.}\\ \noalign{\vskip 5.8pt}
\company{Google Research, Munich Germany}
May 2025 -- Nov 2025 &
  Student Researcher -- 3D Scene Understanding for AR/VR Applications\par
  \bitem{\textit{Advisors}: \href{https://scholar.google.de/citations?user=dfjN3YAAAAAJ&hl=en}{Johanna Wald} \& \href{https://federicotombari.github.io/}{Federico Tombari}}
  \bitemj{\textit{Role}: Research on open-vocabulary 3D scene understanding using forward models, investigating object semantics, affordances \& articulations [1]. Built up a forward model implementation that is used internally across multiple research projects.}\\
\end{cvtable}

\begin{cvtable}{0.831in}
\company{Bosch Center for Artificial Intelligence Stuttgart Germany}
Apr 2022 -- Apr 2025 &
  Bosch Research Doctoral Program -- 3D Scene Understanding for indoor Robotic Applications\par
  \bitem{\textit{Advisors}: \href{https://scholar.google.com/citations?user=k4m1c6EAAAAJ&hl=en&oi=ao}{Mirco Colosi} \& \href{https://scholar.google.com/citations?user=U3KSTwkAAAAJ&hl=en}{Narunas Vaskevicius}}
  \bitemj{\textit{Role}: Research and development on 3D Scene Graphs for indoor 3D scene understanding, achieving over 50\% reduction in labeled sample requirements through self-supervised learning, and introducing open-vocabulary relationship predictions for fine-grained inter-object reasoning. This work resulted in four first-author publications [3], [5], [6], [7].}\\ \noalign{\vskip 2.4pt}
Oct 2021 -- Mar 2022 &
  Industrial Master Thesis -- 6D Pose Estimation\par
  \bitem{\textit{Advisor}: \href{https://scholar.google.com/citations?user=efxPCRAAAAAJ&hl=en&oi=ao}{Fabian Duffhauss}}\\
 &
  \bitemcont{\textit{Role}: Developed a multi-view RGB-D fusion method with a symmetry-aware keypoint voting approach, improving 6D Pose Estimation accuracy by 7.6\% and achieving state-of-the-art results, resulting in a publication [8].}\\
\end{cvtable}

\begin{cvtable}{0.831in}
\company{University of Tübingen, Tübingen Germany}
Sep 2020 -- Oct 2021 &
  Research Assistant -- Real-time high-resolution remote sensing\par
  \bitem{\textit{Advisors}: \href{https://scholar.google.com/citations?hl=en&user=4XOjpZ8AAAAJ&view_op=list_works}{Leon-Amadeus Varga} \& \href{https://uni-tuebingen.de/fakultaeten/mathematisch-naturwissenschaftliche-fakultaet/fachbereiche/informatik/lehrstuehle/kognitive-systeme/the-chair/staff/prof-dr-andreas-zell/}{Prof.\ Andreas Zell}}
  \bitemj{\textit{Role}: Research on encoder-only region proposal techniques combined with hardware-optimized TensorRT and CUDA implementations on embedded GPUs for real-time object detection in high-resolution images. Studied the impact of on-device image pre-processing for enhanced remote sensing object detection accuracy, resulting in a publication [9].}\\
\end{cvtable}

\begin{cvtable}{0.831in}
\company{Bosch Research, Stuttgart Germany}
Apr 2020 -- Oct 2020 &
  Research Intern -- Object-level Semantics for SLAM\par
  \bitem{\textit{Advisor}: \href{https://www.linkedin.com/in/stefan-benz-95940b154/?originalSubdomain=de}{Stefan Benz}}
  \bitemj{\textit{Role}: Integrated an object detection pipeline into a ROS system for enhanced localization and mapping. Led synthetic data generation using Unreal Engine to enable reproducible mapping runs used in the evaluation section of a published paper.}\\
\end{cvtable}

\begin{cvtable}{0.831in}
\company{Bosch Group, Stuttgart Germany}
Oct 2016 -- Mar 2020 &
  Industrial Bachelor Student\par
  \bitemj{\textit{Role}: Alongside my bachelor's in computer science at the university, I worked on scientific projects in different departments at Bosch, applying learned knowledge and collaborating with experienced developers and researchers.}
  \bitem{\textit{Mentors}: Stefan Benz, Hanna Ziesche, Christopher Baker, Stephan Stühmer}
  \bitems{\textit{Nov 2017 -- Feb 2018}: Analyzed false-positive ultrasonic reflections in automotive sensors.}
  \bitems{\textit{Jun 2018 -- Sep 2018}: Implemented a play-back recording feature for ECU failure analysis.}
  \bitems{\textit{Jun 2019 -- Sep 2019}: Implemented a pedestrian simulator for ML-based human behavior prediction.}
  \bitems{\textit{Sep 2019 -- Mar 2020}: Researched the impact of semantic features for improved SLAM accuracy.}\\
\end{cvtable}

% ======================================================== honors & awards ===
\setlength{\entrygap}{3.6pt}% compact sections use tighter gaps
\cvsection[12.2pt]{Honors \& Awards}

\begin{cvtable}[\compactsize]{0.421in}
2024 & Accepted into International Computer Vision Summer School (ICVSS) 2024 for excellent PhDs in CV.\\
\end{cvtable}

\begin{cvtable}[\compactsize]{0.421in}
2021 & 1st place in the AI Chess Competition conducted by the Cognitive Systems Lab of Prof.\ Andreas Zell.\par
       3rd place in the RL Hockey Competition of the MPI Autonomous Learning Group of Prof.\ Georg Martius.\\
\end{cvtable}

\begin{cvtable}[\compactsize]{0.421in}
2020 & Accepted into the Students@Bosch program for students who excelled at Bosch internships.\\
\end{cvtable}

\begin{cvtable}[\compactsize]{0.421in}
2014 & 4th place at the RoboCup World Cup 2014 in Joao Pessoa Brazil in the Rescue Junior competition\\
\end{cvtable}

\begin{cvtable}[\compactsize]{0.421in}
2013 & German Champion at the RoboCup German Open 2013 in the Rescue Junior competition.\par
       World Champion at the RoboCup World Cup 2013 in Eindhoven in the Rescue Junior competition.\\
\end{cvtable}

% ========================================================== publications ====
\cvsection[10.6pt]{Publications}
\vspace{-4pt}
{\fontsize{8}{9.6}\selectfont For a complete list of all publications, including recent pre-preprints see: \textit{\href{http://kochsebastian.com/publications}{kochsebastian.com/publications}}\par}
\vspace{4.4pt}

\begin{cvtable}[\compactsize]{0.441in}
2026 &
  \pubitem{[1] \textbf{S Koch}, J Wald, H Matsuki, P Hermosilla, T Ropinski, F Tombari, \textit{``Unified Semantic Transformer for\newline 3D Scene Understanding'',} in Transactions on Machine Learning Research (\textbf{TMLR}), 2026.}\\ \noalign{\vskip 2.6pt}
2025 &
  \pubitem{[2] L Weijler, \textbf{S Koch}, F Poiesi, T Ropinski, P Hermosilla, \textit{``OpenHype: Hyperbolic Embeddings for Hierarchical Open-Vocabulary Radiance Fields'',} in Advances in Neural Information Processing Systems (\textbf{NeurIPS}), 2025.}\\ \noalign{\vskip 2.6pt}
 &
  \pubitem{[3] \textbf{S Koch}, J Wald, M Colosi, N Vaskevicius, P Hermosilla, F Tombari, T Ropinski, \textit{``RelationField: Relate Anything in Radiance Fields'',} in Conference on Computer Vision and Pattern Recognition (\textbf{CVPR}),\newline 2025.}\\ \noalign{\vskip 2.6pt}
 &
  \pubitem{[4] Y Liu, L Palmieri, \textbf{S Koch}, I Georgievski, M Aiello, \textit{``DELTA: Decomposed Efficient Long-Term Robot Task Planning using Large Language Models''}, in IEEE International Conference on Robotics and Automation \textbf{(ICRA)}, 2025.}\\
\end{cvtable}

\begin{cvtable}[\compactsize]{0.441in}
2024 &
  \pubitem{[5] \textbf{S Koch}, N Vaskevicius, M Colosi, P Hermosilla, T Ropinski, \textit{``Open3DSG: Open-vocabulary 3D Scene Graphs from Point Clouds with Queryable Objects and Open-Set Relationships'',} in Conference on Computer Vision and Pattern Recognition (\textbf{CVPR}), 2024.}\\ \noalign{\vskip 2.6pt}
 &
  \pubitem{[6] \textbf{S Koch}, P Hermosilla, N Vaskevicius, M Colosi, T Ropinski, \textit{``Lang3DSG: Language-based contrastive pre-training for 3D Scene Graph prediction''}, in International Conference on 3D Vision (\textbf{3DV}), 2024.}\\ \noalign{\vskip 2.6pt}
 &
  \pubitem{[7] \textbf{S Koch}, P Hermosilla, N Vaskevicius, M Colosi, T Ropinski, \textit{``SGRec3D: Self-Supervised 3D Scene Graph Learning via Object-Level Scene Reconstruction'',} in Winter Conference on Applications of Computer Vision (\textbf{WACV}), 2024.}\\
\end{cvtable}

\begin{cvtable}[\compactsize]{0.441in}
2023 &
  \pubitem{[8] F Duffhauss, \textbf{S Koch}, H Ziesche, NA Vien, G Neumann, \textit{``SyMFM6D: Symmetry-aware Multi-directional Fusion for Multi-View 6D Object Pose Estimation'',} in Robotics and Automation Letters (\textbf{RA-L}), 2023.}\\
\end{cvtable}

\begin{cvtable}[\compactsize]{0.441in}
2022 &
  \pubitem{[9] LA Varga, \textbf{S Koch}, A Zell, \textit{``Comprehensive Analysis of the Object Detection Pipeline on UAVs''}, in \textbf{Remote Sensing}, 2022.}\\
\end{cvtable}

% ========================================================= invited talks ====
\cvsection[10.4pt]{Invited Talks}

\begin{cvtable}[\compactsize]{0.441in}
2025 &
  \textbf{Google Research:} \textit{Open-Vocabulary Relationship Reasoning in 3D}, Munich 05.05.2025.\par
  \vspace{3.5pt}
  \textbf{Huawai Munich Research Center}: \textit{Language-driven Scene Understanding with 3D Scene Graphs},\newline Munich 06.02.2025.\\
\end{cvtable}

\begin{cvtable}[\compactsize]{0.441in}
2024 &
  \textbf{RWTH Aachen.\ Computer Vision Group, Prof.\ Bastian Leibe}: \textit{3D Scene Understanding with label-efficient and open-vocabulary 3D Scene Graphs}, Remote 27.02.2024.\par
  \vspace{3.5pt}
  \textbf{Bosch Research US, Sunnyvale.} \textit{3D Scene Understanding with label-efficient and open-vocabulary 3D Scene Graphs}, Remote 07.02.2024.\\
\end{cvtable}

% =============================================================== service ====
\cvsection[9pt]{Service}

\begingroup\fontsize{10.5}{13.6}\selectfont
{\hangindent=1.24in \hangafter=1 \noindent
\textbf{Reviewing}: \underline{Conferences}: CVPR 2024--2026, ICCV 2023--2025, ECCV 2026, NeurIPS 2024--2025, ICML 2026, ICRA 2025, ROS 2024--2025 --- \underline{Journals}: RA-L 2024--2026\newline \underline{Recognition}: Outstanding Reviewer: CVPR 2025, CVPR 2026\par}
\textbf{Volunteering}: I work as a volunteer and referee at RoboCup Junior events on a national and international level.\par
\endgroup

% ======================================================== qualifications ====
\cvsection[14.5pt]{Qualifications}

\begingroup\fontsize{10.5}{13.6}\selectfont
\textbf{Programming}: Python, C/C++, MATLAB, Bash\par
\textbf{Libraries}: PyTorch, JAX, NumPy, Lightning, FLAX, Matplotlib, Open3D, OpenCV, TensorRT\par
\textbf{Software}: Git, \LaTeX, ROS, GNU/Linux, Blender, Office\par
\textbf{Languages:} German (native), English (fluent)\par
\endgroup

% ============================================================ references ====
\cvsection[13pt]{References}
\vspace{3pt}

\newcommand{\reference}[3]{\noindent\textbf{#1.} #2\hfill\href{mailto:#3}{#3}\par}
\begingroup\fontsize{10.5}{13.6}\selectfont
\reference{Prof.\ Timo Ropinski}{Full Professor, Ulm University.}{timo.ropinski@uni-ulm.de}
\reference{Dr.\ Federico Tombari}{Director of Research, Google Zurich.}{tombari@google.com}
\reference{Prof.\ Pedro Hermosilla}{Assistant Professor, TU Wien.}{phermosilla@cvl.tuwien.ac.at}
\reference{Dr.\ Narunas Vaskevicius}{Lead Research Scientist, Bosch Research.}{narunas.vaskevicius@bosch.com}
\reference{Prof.\ Gerard Pons-Moll}{Full Professor, University of Tuebingen.}{gerard.pons-moll@uni-tuebingen.de}
\endgroup

\end{document}
